\section*{Chapitre \ref{ch:b3:photochemistry} --- Photochimie}

\begin{solution}{exo:b3:photochemistry:1}
$E = hc/\lambda = \qty{5.442e-19}{J}$~; $0.100/\num{5.442e-19} = \num{1.84e17}$ photons par seconde, soit
\qty{3.05e-7}{einstein.s^{-1}}.
\end{solution}

\begin{solution}{exo:b3:photochemistry:2}
Absorbé~: $\num{1.00e-7} \times 600 = \qty{6.00e-5}{einstein}$~; $\Phi = \num{2.0e-5}/\num{6.0e-5} = 0.33$.
\end{solution}

\begin{solution}{exo:b3:photochemistry:3}
Bien plus de molécules réagissent que de photons ne sont absorbés~:
chaque \ce{Cl2} photolysé amorce une
chaîne, $\ce{Cl + H2 -> HCl + H}$ et $\ce{H + Cl2 -> HCl + Cl}$, qui
accomplit des milliers de tours avant la
rupture.
\end{solution}

\begin{solution}{exo:b3:photochemistry:4}
$D_0(\ce{O2}) = 2 \times 246.79 = \qty{493.58}{kJ/mol}$~; $\lambda = hcN_A/D_0 = 0.119627/493\,580 =
\qty{242.4}{nm}$. $\ce{O3 -> O2 + O}$: $246.79 - 145.35 = \qty{101.44}{kJ/mol}$, $\lambda = \qty{1179}{nm}$.
\end{solution}

\begin{solution}{exo:b3:photochemistry:5}
$[Z]/[E] = 22\,000 \times 0.11/(1500 \times 0.40) = 4.0$~: 80\,\% de $Z$, 20\,\% de $E$.
\end{solution}

\begin{solution}{exo:b3:photochemistry:6}
Hexan-2-one~: le biradical-1,4 se coupe en énol de l'acétone (donc en
acétone) et en
propène, ou se referme en 1,2-diméthylcyclobutan-1-ol. Dans la
pentan-2-one, le carbone $\gamma$
est le \ce{CH3} terminal~: la coupure de la liaison $\alpha$--$\beta$ libère $\ce{CH2=CH2}$.
\end{solution}

\begin{solution}{exo:b3:photochemistry:7}
$K = \exp[(E_{\mathrm T}(\mathrm D) - 250)/RT]$ avec $RT = \qty{2.478}{kJ/mol}$~: 289 donne $\eu^{15.7} = \num{7e6}$~;
253 donne $\eu^{1.2} = 3.4$~; 236 donne $\eu^{-5.6} = \num{3.5e-3}$. Seul
le premier transfère presque à chaque rencontre.
\end{solution}

\begin{solution}{exo:b3:photochemistry:8}
$k_2 = \num{6.0e-34} \times (230/300)^{-2.6} = \qty{1.20e-33}{cm^6.s^{-1}}$~; $k_2[\mathrm M][\ce{O2}] = \num{1.20e-33} \times
\num{4.0e17} \times \num{8.4e16} = \qty{40}{s^{-1}}$; $[\ce{O}]/[\ce{O3}] = \num{1.0e-3}/40 = \num{2.5e-5}$.
\end{solution}

\begin{solution}{exo:b3:photochemistry:9}
$k(\ce{Cl + O3}) = \num{2.8e-11}\eu^{-250/230} = \num{9.4e-12}$~; $k(\ce{Cl + CH4}) = \num{6.6e-12}\eu^{-1240/230} = \num{3.0e-14}$
\unit{cm^3.s^{-1}}. Vitesses de pseudo-ordre 1~: 54 et \qty{0.012}{s^{-1}}~; probabilité $54/54.01 = 0.9998$.
\end{solution}

\begin{solution}{exo:b3:photochemistry:10}
$k = \num{2.07e-12}\eu^{-1400/298.15} = \qty{1.89e-14}{cm^3.s^{-1}}$~; $[\ce{O3}] = \num{8.0e-3} \times 1.5/\num{1.89e-14} =
\qty{6.3e11}{cm^{-3}}$. L'air sous \qty{1}{bar} et à \qty{298}{K}
contient $p/k_BT = \qty{2.43e19}{cm^{-3}}$~: \qty{26}{ppb}. La nuit,
\ce{NO2} n'est plus photolysé et NO détruit l'ozone jusqu'à épuisement de
l'un des deux.
\end{solution}

\begin{solution}{exo:b3:photochemistry:11}
$E = hc/\lambda = \qty{6.347e-19}{J}$~; $0.050/\num{6.347e-19} = \num{7.88e16}$ photons par seconde $= \qty{1.31e-7}{einstein.s^{-1}}$.
Absorbé~: $\times(1 - 10^{-0.5}) = \times0.684$, \qty{8.95e-8}{einstein.s^{-1}}~; vitesse $0.25 \times \num{8.95e-8} =
\qty{2.24e-8}{mol.s^{-1}}$, soit \qty{7.5e-6}{mol.L^{-1}.s^{-1}} dans \qty{3.0}{mL}.
\end{solution}

\begin{solution}{exo:b3:photochemistry:12}
$I_0 = \num{3.0e-6}/(1.25 \times 60) = \qty{4.0e-8}{einstein.s^{-1}}$. À
faible conversion, l'actinomètre absorbe encore toute la lumière, et les
produits colorés ne la filtrent pas.
\end{solution}

\begin{solution}{pb:b3:photochemistry:1}
\textbf{1.} $2 \times 246.79 = \qty{493.58}{kJ/mol}$.
\textbf{2.} $\lambda = 0.119627/493\,580 = \qty{242.4}{nm}$.
\textbf{3.} $246.79 - 145.35 = \qty{101.44}{kJ/mol}$~; \qty{1179}{nm}.
\textbf{4.} L'absorption visible de l'ozone est faible~; sa forte
absorption ultraviolette (d'environ 200 à \qty{310}{nm}) élimine le
rayonnement qui endommage l'ADN et la peau.
\textbf{5.} La lumière de longueur d'onde inférieure à \qty{242}{nm} est
absorbée par \ce{O2} lui-même au cours de sa descente~: peu en pénètre
au-dessous de la haute stratosphère.
\textbf{6.} $\ce{O2 ->[h\nu] 2 O}$~; $\mathrm{O + O_2 + M \to O_3 + M}$~; $\ce{O3 ->[h\nu] O2 + O}$~; $\ce{O + O3 -> 2 O2}$.
\textbf{7.} $k_2 = \qty{1.20e-33}{cm^6.s^{-1}}$~; $k_2[\mathrm M] = \qty{4.79e-16}{cm^3.s^{-1}}$.
\textbf{8.} $k_4 = \num{8.0e-12}\eu^{-2060/230} = \qty{1.03e-15}{cm^3.s^{-1}}$.
\textbf{9.} $\num{1.0e-3}/(\num{4.79e-16} \times \num{8.4e16}) = \num{2.49e-5}$.
\textbf{10.} Bilan de l'oxygène impair $2j_1[\ce{O2}] = 2k_4[\ce{O}][\ce{O3}]$ avec $[\ce{O}] = j_3[\ce{O3}]/(k_2[\mathrm M][\ce{O2}])$
(voir le théorème).
\textbf{11.} $[\ce{O3}] = \num{8.4e16}\sqrt{\num{1.0e-11} \times \num{4.79e-16}/(\num{1.0e-3} \times \num{1.03e-15})} =
\qty{5.7e12}{cm^{-3}}$, soit une fraction molaire de $\num{1.4e-5}$ (\qty{14}{ppm}).
\textbf{12.} $\num{2.49e-5} \times \num{5.7e12} = \qty{1.4e8}{cm^{-3}}$.
\textbf{13.} Des cycles catalytiques (oxydes d'azote, oxydes d'hydrogène,
chlore) ajoutent des voies de destruction de l'oxygène impair.
\textbf{14.} $\ce{Cl + O3 -> ClO + O2}$, $\ce{ClO + O -> Cl + O2}$~; bilan $\ce{O + O3 -> 2 O2}$.
\textbf{15.} $1/(\num{9.44e-12} \times \num{5.72e12}) = \qty{0.019}{s}$.
\textbf{16.} $1/(\num{4.033e-11} \times \num{1.423e8}) = \qty{174}{s}$.
\textbf{17.} $174/0.0185 = \num{9.4e3}$.
\textbf{18.} Pratiquement tout le chlore est sous forme ClO~: destruction
$2 \times \num{5.74e-3} \times \num{1.0e7} = \qty{1.1e5}{cm^{-3}.s^{-1}}$,
contre $2k_4[\ce{O}][\ce{O3}] = \qty{1.7e6}{cm^{-3}.s^{-1}}$ pour l'étape
de Chapman~: le cycle ajoute environ 7\,\%.
\textbf{19.} $\ce{ClO + O}$, l'étape lente (\qty{174}{s}).
\textbf{20.} $53.8/(53.8 + 0.012) = 0.99978$.
\textbf{21.} $\num{5.74e-3}/(\num{5.74e-3} + \num{1.41e-13} \times \num{1.0e9}) = 0.976$.
\textbf{22.} $p = 0.976$.
\textbf{23.} $0.976/0.024 = 40$.
\textbf{24.} Sur les nuages, $\ce{HCl + ClONO2 -> Cl2 + HNO3}$ libère le
chlore, et l'acide nitrique reste dans les particules~: sans \ce{NO2},
$p_2$ tend vers 1 et la chaîne devient très longue~; un cycle passant par
le dimère de ClO détruit l'ozone sans atomes O.
\textbf{25.} \textbf{Environ 40 molécules d'ozone par atome de chlore avant
sa capture} (dans ce modèle à \qty{30}{km})~; chaque atome est libéré de
ses réservoirs à maintes reprises.
\end{solution}
